% ============================================================================
%  Entregable 2. Análisis económico-financiero (Formulario E-21, FEP01 p.71).
%  Apartados 2.1 a 2.7 del formulario. Las cifras salen de
%  economico/datos/partidas.csv: aquí no se escribe ningún monto a mano.
% ============================================================================

\begin{resumenApertura}
\marcador{Qué justifica este análisis, en cuatro o cinco líneas, con el costo
total frente al valor ofertado}
\begin{recibe}
  \item \marcador{Qué recibe el mandante}
\end{recibe}
\end{resumenApertura}

\section{Análisis comparativo de costos frente a precio}\label{sec:eco-costo-precio}

\marcador{Párrafo de contexto: qué muestra la \tabref{eco-maestro}, el cuadro
maestro de costos por categoría}

\tablaMontos{Cuadro maestro de costos}{eco-maestro}%
  {costo-directo-implementacion, costo-directo-operacion, costo-indirecto}

\marcador{Párrafo de contexto: qué muestra la \tabref{eco-precio}, el precio de
venta con el margen aplicado a cada componente}

\begin{tabla}{Precio de venta y margen por componente}{eco-precio}{X R{30mm} R{22mm} R{30mm}}
Componente & Costo & Margen & Precio de venta \\ \midrule
\marcador{componente} & \marcador{partida de costo} & \marcador{porcentaje} & \marcador{partida de precio} \\
\end{tabla}

\marcador{Conclusión que se saca de las dos tablas}

\section{Evaluación financiera}\label{sec:eco-evaluacion}

\marcador{Párrafo de contexto: qué muestra la \tabref{eco-indicadores}. VAN con
la tasa de descuento y su justificación, flujo de caja proyectado, TIR,
sensibilidad y punto de equilibrio}

\begin{tabla}{Indicadores de la evaluación financiera}{eco-indicadores}{L{60mm} X}
Indicador & Valor \\ \midrule
Horizonte de evaluación & 56 meses, \form{E-24}{73} \\
Tasa de descuento & \marcador{tasa y justificación} \\
Valor actual neto (VAN) & \montoPartida{van} \\
Tasa interna de retorno (TIR) & \marcador{porcentaje} \\
Punto de equilibrio & \marcador{mes} \\
\end{tabla}

\marcador{Análisis de sensibilidad: qué variables se mueven y cuánto cambia el VAN}

\section{Curva S}\label{sec:eco-curva}

\marcador{Costos acumulados, ingresos acumulados y análisis del flujo de caja
mensual. La figura se cita con \textbackslash figref antes de aparecer}

\section{Estrategia de adquisiciones}\label{sec:eco-estrategia}

\marcador{Matriz de adquisiciones principales, cronograma de compras alineado con
el plan de proyecto, estrategia de negociación, contratos críticos, plan de
importaciones si aplica y análisis de hacer o comprar}

\section{Modelo de costos operacionales}\label{sec:eco-operacion}

\marcador{Recursos humanos, infraestructura y alojamiento, licenciamiento y
suscripciones, conectividad, energía e instalaciones, seguros y garantías, y
optimizaciones previstas en el tiempo}

\subsection{Costo y tarifa por perfil}\label{sec:eco-tarifas}

\marcador{Párrafo de contexto: qué muestra la \tabref{tarifas-e26}. El costo y la
tarifa ofertados van en economico/datos/tarifas.csv y deben caer en los rangos
del Formulario E-26}

\tablaTarifas

\section{Estructura de mantención y soporte}\label{sec:eco-mantencion}

\marcador{Modelo de costeo, recursos dedicados frente a compartidos, mantención
preventiva programada, provisión para mantención correctiva, presupuesto de
mejora continua y costos de actualización tecnológica}

\section{Valorización de las cinco innovaciones}\label{sec:eco-innovaciones}

\marcador{Párrafo de contexto: qué muestra la \tabref{eco-innovaciones}, la
inversión, el costo operacional incremental y el beneficio esperado de cada
innovación, reflejados en el flujo de caja}

\tablaMontos*[fuente=condensada]{Valorización de las cinco innovaciones}{eco-innovaciones}%
  {innovacion-1-inversion, innovacion-1-costo, innovacion-1-beneficio,
   innovacion-2-inversion, innovacion-2-costo, innovacion-2-beneficio,
   innovacion-3-inversion, innovacion-3-costo, innovacion-3-beneficio,
   innovacion-4-inversion, innovacion-4-costo, innovacion-4-beneficio,
   innovacion-5-inversion, innovacion-5-costo, innovacion-5-beneficio}

\marcador{Conclusión que se saca de la tabla}
